\section{Introduction}
\label{sec:intro}

Deep-research agents---systems that plan, search the web, and synthesize long reports with
citations---are now evaluated almost entirely on \emph{correctness}: whether their claims are
factual and whether citations support them \citep{du2025deepresearchbench,liveresearchbench2025,
mialon2023gaia,wei2025browsecomp}. Correctness is necessary but not sufficient for the use case
these agents are marketed for: helping researchers, analysts, and inventors find ideas that are
not already obvious. That capacity is \emph{creativity}, and it is essentially unmeasured for
this class of system.

We introduce \textbf{CREST} (CREativity of reSearch agenTs), the first benchmark that scores
deep-research outputs for creativity using the standard construct from the creativity literature:
an output is creative only if it is simultaneously \emph{New}, \emph{Useful}, and
\emph{Surprising} \citep{runco2012standard,boden2004creative,amabile1982cat}. CREST organizes
questions by \emph{task type}---inventing under constraints, counterfactual reasoning,
cross-domain synthesis, meta-creativity, and open-ended generation---because our central finding
is that creativity in these agents is not a single scalar but \emph{depends on the task}.

\paragraph{The finding.} Running CREST on an exploration-based research agent (anonymized here;
its architecture is described in a companion preprint we cite in the third person
\citep{caesar2026}) alongside frontier and open-weight competitors, we find the agent is more
creative on \emph{cross-domain} and \emph{counterfactual} tasks, roughly at parity on
\emph{meta-creativity}, and \emph{not} more creative on \emph{constrained} and \emph{open-ended}
tasks. Ablations tie each win to a specific mechanism: turning off one component degrades one
task type and not the others. The wins survive a human study and coexist with strong
citation-faithfulness, so the creative reports are not confident fabrication.

\paragraph{Why the losses matter.} We deliberately foreground where the agent \emph{loses}. A
method that improved every task type would be indistinguishable from a length or judge-bias
artifact. Winning only where the mechanism predicts---and losing where it does not, on a
benchmark whose own author's system does not top every category---is the strongest available
evidence that the effect is real and mechanism-driven rather than an artifact of a self-serving
test \citep{singh2025leaderboard}. The hero of this paper is \textbf{CREST + the task-dependence
finding}, not ``our agent wins.''

\paragraph{Contributions.}
\begin{enumerate}
  \item \textbf{Finding.} Creativity in deep-research agents is task-dependent, with ablation
    evidence that specific architectural components cause specific gains (\S\ref{sec:experiments}).
  \item \textbf{Benchmark.} CREST---the first creativity benchmark for deep-research agents:
    New/Useful/Surprising scoring by task type (fifty questions, ten per type), a human-validated
    leave-one-family-out median judge panel, a frozen shared retrieval index for reproducibility,
    and a full public release
    (\S\ref{sec:crest}).
  \item \textbf{Evidence of creative-and-grounded synthesis.} The same reports that score high on
    creativity rank competitively on two established citation-faithfulness benchmarks
    (\S\ref{sec:grounding}).
\end{enumerate}

We state the scope honestly (\S\ref{sec:limitations}): the mechanism evidence is currently
strongest for one agent, we test the task-dependence pattern across several agents, and we
report the two task types where our judge and human experts disagree.
